% TEMPLATE-MILC-v3.tex - plantilla maestra UNICA de las guias MILC v3.
%
% La usan las cuatro familias de guias, cada una con su driver:
%   · Grados 6-11          -> scripts/build-guias-g11.py
%   · Bebras               -> scripts/build-guias-bebras.py
%   · Semillero            -> scripts/build-guias-semillero.py
%   · Territorio interior  -> scripts/build-guias-territorio-interior.py
%
% Antes habia tres copias casi identicas (~400 lineas iguales, 6 distintas):
% cada mejora habia que aplicarla tres veces, y en la practica no se hacia
% (el motor de recursos vivia solo en dos de ellas). Lo que varia entre
% familias esta parametrizado en 6 placeholders que cada driver llena
% (se nombran aqui SIN su sintaxis real: el builder reemplaza por texto plano,
% tambien dentro de los comentarios, y un valor multilinea rompe el preambulo):
%
%   HEADER_IZQ        encabezado izquierdo de pagina
%   HEADER_DER        encabezado derecho de pagina
%   PORTADA_ETIQUETA  rotulo sobre el numero grande (GRADO/MOMENTO/MODULO)
%   PORTADA_SERIE     linea de serie de la portada
%   PORTADA_ID        identificador de la guia en la portada
%   PORTADA_CIRCULO   el circulito de grado (°), o vacio si no aplica
%   PDF_TITULO        titulo en texto plano (sin \textbf ni comillas TeX) para
%                     los metadatos del PDF (/Title); lo produce lib_xelatex.texto_pdf
%
% Composicion (auditoria 2026-09-03): las bandas de estacion piden espacio
% (\Needspace*) y prohiben el salto justo despues (\nopagebreak), asi no quedan
% huerfanas al pie; las cajas largas (softbox, drawbox, infoband, puentebox) son
% `breakable`, asi una caja de media pagina ya no arrastra todo a la siguiente
% dejando la anterior al 40 %. Todo texto sobre mostaza o verde va en milcNegro
% (blanco sobre #E5B400 da 1,93:1 y sobre #58A923 2,95:1; ninguno pasa WCAG AA).
% El resto de claves las documenta content/guias/_SCHEMA.md. El script valida
% que no queden placeholders sin reemplazar antes de escribir el archivo final.

% Etiquetado PDF (latex-lab): /Lang, estructura y texto alternativo de las
% figuras, para lectores de pantalla. Debe ir ANTES de \documentclass.
\DocumentMetadata{lang=es-CO,tagging=on}
\documentclass[11pt,letterpaper]{article}
% headheight=24pt: la cabecera derecha lleva dos lineas (identidad de la guia
% y estacion actual). headsep se acorta en la misma medida para que el bloque
% de texto quede exactamente donde estaba con headheight=12pt.
\usepackage[margin=2.0cm,headheight=24pt,headsep=13pt]{geometry}
\usepackage[table]{xcolor}
\usepackage{fontspec}
\usepackage[spanish]{babel}
\usepackage{tikz}
% Librerías TikZ para los diagramas de las guías (bloque `recursos:`):
%  · babel  → babel en español vuelve activos caracteres como «>», y TikZ los usa
%             en sus opciones (>=stealth, ->). Sin esta librería, cualquier
%             diagrama con flechas revienta con «\language@active@arg> has an
%             extra }». Debe ir ANTES de que se dibuje nada.
%  · positioning        → sintaxis «below left=2pt and 3pt».
%  · decorations.pathreplacing → llaves ({brace}) para señalar rangos.
\usetikzlibrary{babel,positioning,decorations.pathreplacing}
\usepackage{graphicx}
% Circuitos: el trabajo de electrónica se hace hoy con micro:bit y sus sensores
% integrados (MakeCode, sin cableado), así que no se necesitan esquemáticos.
% Si algún día entran montajes con protoboard, basta descomentar esta línea:
% los diagramas .tex/.tikz declarados en `recursos:` ya se incrustan solos.
% \usepackage{circuitikz}
\usepackage[most]{tcolorbox}
\usepackage{tabularx}
\usepackage{array}
\usepackage{fancyhdr}
\usepackage{titlesec}
\usepackage[document]{ragged2e}  % bandera a la derecha en todo el cuerpo
\usepackage{enumitem}
\usepackage{microtype}
\usepackage{etoolbox}
\usepackage{needspace}
% hyperref al final del preambulo: metadatos del PDF (titulo, autor, idioma,
% familia), marcadores por estacion (\pdfbookmark) y enlaces sin recuadro.
\usepackage[hidelinks,
  pdftitle={El que llegó de otra parte: hasta la ciudad tuvo que mudarse},
  pdfauthor={PhD. Álvaro Cárdenas Orozco},
  pdfsubject={Territorio interior · Educación socioemocional · Grado 9° · Momento 7 · MILC v3},
  pdflang={es-CO},
  bookmarksopen=true,bookmarksnumbered=false]{hyperref}

% Helvetica Neue no trae la flecha U+2192, asi que cada «→» escrito literalmente
% en el YAML se compilaba a NADA, en silencio (462 flechas en 61 guias: rutas del
% tipo «paleta → tipografia → lockup» salian rotas). Mapeamos el caracter a su
% version matematica, que si tiene glifo. Asi el autor puede seguir escribiendo
% «→» comodamente en el YAML.
\usepackage{newunicodechar}
\newunicodechar{→}{\ensuremath{\rightarrow}}
% Mismo problema con los simbolos de marca y vinetas: el «✓» de «una palomita ✓»
% salia como cuadrito vacio en el PDF de todas las guias que lo usaban.
\usepackage{pifont}
\newunicodechar{✓}{\ding{51}}
\newunicodechar{✗}{\ding{55}}
\newunicodechar{✔}{\ding{52}}
\newunicodechar{•}{\textbullet}
\newunicodechar{≥}{\ensuremath{\geq}}
\newunicodechar{≤}{\ensuremath{\leq}}

\setmainfont{Helvetica Neue}
\setsansfont{Helvetica Neue}
\setlength{\parindent}{0pt}
\setlength{\parskip}{6pt}
% Sin particion de palabras: en 6.º-7.º una palabra cortada por guion al final
% de linea («Comuni-cacion») frena la lectura mucho mas que una linea corta.
\hyphenpenalty=10000
\exhyphenpenalty=10000
\pretolerance=2000
\tolerance=3000
\setlength{\emergencystretch}{3em}
\linespread{1.10}
\renewcommand{\arraystretch}{1.25}
\setlist{leftmargin=5mm,itemsep=1mm,topsep=1mm}
\newcolumntype{Y}{>{\RaggedRight\arraybackslash}X}

\definecolor{milcMagenta}{HTML}{E6007E}
\definecolor{milcVino}{HTML}{5A0038}
\definecolor{milcTurquesa}{HTML}{0093A5}
\definecolor{milcVerde}{HTML}{58A923}
\definecolor{milcMostaza}{HTML}{E5B400}
\definecolor{milcOcre}{HTML}{B86600}
\definecolor{milcNegro}{HTML}{241816}
\definecolor{milcGris}{HTML}{F7F7F4}
\definecolor{milcLinea}{HTML}{E8E2DD}
\definecolor{milcGradePrimary}{HTML}{0D9488}
\definecolor{milcGradeDark}{HTML}{115E59}
\definecolor{milcGradeSoft}{HTML}{CCFBF1}
\definecolor{milcGradeAccent}{HTML}{CFE7FF}
\definecolor{milcGradeText}{HTML}{FFFFFF}
\color{milcNegro}

\pagestyle{fancy}
\fancyhf{}
% Cabecera de dos lineas: identidad de la guia arriba y, a la derecha y en
% gris, la estacion en curso (\markright desde \milcestacion). Antes de la
% primera estacion la segunda linea va vacia; el \strut mantiene la altura
% para que la linea de arriba no baile entre paginas.
\lhead{\small Territorio interior · Educación socioemocional\\[-1pt]{\footnotesize\strut}}
\rhead{\small Grado 9° · Momento 7 · MILC v3\\[-1pt]{\footnotesize\strut\color{milcNegro!65}\rightmark}}
\cfoot{\small\thepage}
\renewcommand{\headrulewidth}{0pt}

% Sobre mostaza (#E5B400) y verde (#58A923) el texto va en milcNegro; sobre el
% resto de la paleta, en blanco. Se usa dentro de fonttitle/fontupper, donde un
% \color posterior manda sobre coltitle.
\newcommand{\milctintasobre}[1]{%
  \ifboolexpr{test{\ifstrequal{#1}{milcMostaza}} or test{\ifstrequal{#1}{milcVerde}}}%
    {\color{milcNegro}}{\color{white}}}

\newtcolorbox{titlebox}[1]{
  enhanced,colback=#1,colframe=#1,arc=7mm,boxrule=0pt,
  left=7mm,right=7mm,top=3mm,bottom=3mm,
  before skip=8pt,after skip=8pt,
  fontupper=\sffamily\bfseries\Large\color{white}
}

\newtcolorbox{softbox}[3]{
  enhanced,breakable,pad at break=3mm,
  colback=#2,colframe=#1,borderline west={4pt}{0pt}{#1},
  arc=2mm,boxrule=.4pt,left=4mm,right=4mm,top=3mm,bottom=3mm,
  before skip=6pt,after skip=8pt,title={#3},coltitle=white,
  colbacktitle=#1,fonttitle=\sffamily\bfseries\milctintasobre{#1},fontupper=\RaggedRight,
  attach boxed title to top left={xshift=3mm,yshift=-2mm},
  boxed title style={arc=2mm, boxrule=0pt}
}

\newtcolorbox{drawbox}[2]{
  enhanced,breakable,pad at break=4mm,
  colback=white,colframe=#1,arc=4mm,boxrule=1pt,
  left=5mm,right=5mm,top=4mm,bottom=4mm,before skip=8pt,after skip=8pt,
  title={#2},coltitle=white,colbacktitle=#1,fonttitle=\sffamily\bfseries\milctintasobre{#1},
  fontupper=\RaggedRight,
  attach boxed title to top left={xshift=4mm,yshift=-2mm},
  boxed title style={arc=3mm, boxrule=0pt}
}

\newtcolorbox{infoband}[1]{
  enhanced,breakable,pad at break=2mm,colback=#1!8,colframe=#1!50,arc=2mm,boxrule=.5pt,
  left=4mm,right=4mm,top=2mm,bottom=2mm,before skip=4pt,after skip=4pt,
  fontupper=\small\RaggedRight
}

% Puente narrativo entre secciones (contrato editorial Conéctate).
% Da continuidad: "Acabas de X. Ahora vas a Y porque Z."
% Visualmente: banda izquierda en color del grado, texto en itálica pequeña.
\newtcolorbox{puentebox}{
  enhanced,breakable,pad at break=2mm,colback=milcGradeSoft,colframe=milcGradeDark,
  borderline west={3pt}{0pt}{milcGradeDark},
  arc=0mm,boxrule=0pt,left=5mm,right=4mm,top=2.5mm,bottom=2.5mm,
  before skip=4pt,after skip=8pt,
  fontupper=\itshape\small\color{milcNegro}\RaggedRight
}
% La flecha va en modo matematico a proposito: Helvetica Neue NO tiene el glifo
% U+2192, asi que \textrightarrow se compilaba a NADA («Missing character: There
% is no → in font Helvetica Neue») y el puente salia sin flecha. En modo
% matematico la flecha viene de la fuente matematica y si se dibuja.
\newcommand{\puente}[1]{%
  \begin{puentebox}%
    \textcolor{milcGradeDark}{$\rightarrow$}\hspace{2mm}#1%
  \end{puentebox}%
}

\newcommand{\linead}[1]{\noindent\color{milcLinea}\rule{#1}{0.5pt}\color{milcNegro}}
\newcommand{\respuesta}[1]{\par\vspace{2mm}\foreach \n in {1,...,#1}{\noindent\color{milcLinea}\rule{\linewidth}{0.5pt}\par\vspace{5mm}}\color{milcNegro}}
\newcommand{\checkbox}{\textcolor{milcGradeDark}{\fbox{\phantom{x}}}\hspace{5pt}}

% ─── Listas aireadas para las guías socioemocionales ────────────────────
% Componer las verificaciones como «\checkbox texto\\» deja el texto que salta
% de línea debajo del cuadrito y apelmaza el bloque. Estos entornos alinean la
% continuación y airean los ítems. Los usa Territorio Interior; cualquier
% familia puede hacerlo.
\newlist{checklista}{itemize}{1}
\setlist[checklista]{
  label=\textcolor{milcGradeDark}{\fbox{\phantom{x}}},
  leftmargin=9mm, labelsep=3.2mm, itemsep=2.4mm,
  topsep=2.5mm, parsep=0pt, partopsep=0pt, before=\RaggedRight
}
\newlist{pasoslista}{enumerate}{1}
\setlist[pasoslista]{
  label=\textcolor{milcGradeDark}{\sffamily\bfseries\arabic*.},
  leftmargin=8mm, labelsep=3mm, itemsep=2.8mm,
  topsep=1mm, parsep=0pt, partopsep=0pt, before=\RaggedRight
}

\newcommand{\phasecard}[4]{
\begin{tcolorbox}[enhanced,colback=#3,colframe=#2,arc=3mm,boxrule=.5pt,height=3.05cm,valign=center,halign=center,left=1.5mm,right=1.5mm,top=2mm,bottom=2mm]
{\sffamily\bfseries\fontsize{22}{24}\selectfont\textcolor{#2}{#1}}\par
{\sffamily\bfseries\small #4}
\end{tcolorbox}
}

% ── Piezas del rediseno 2026 (legibilidad 6.º-11.º) ───────────────────
% Banda de estacion: reemplaza el titlebox macizo. Titulo a la izquierda,
% dato practico (tiempo/modalidad) a la derecha, en una sola linea.
% Nunca huerfana: pide 9 lineas libres (o salta de pagina rellenando con
% \vfill) y prohibe el salto justo despues de la banda. Nueve y no seis:
% la caja partible que sigue necesita ~80pt (titulo adosado, relleno y dos
% lineas) para arrancar en la misma pagina; con menos, tcolorbox la manda
% entera a la siguiente y la banda se queda sola. El penalty 10000 va
% en `after`, ANTES del espacio posterior: un `after skip` (aunque sea 0pt)
% es un \vskip tras la caja, y ese glue es un punto de corte legal que un
% \nopagebreak escrito despues de \end{tcolorbox} ya no cierra. Cada banda
% deja marcador PDF y actualiza la cabecera derecha.
\newcounter{milcestacion}
\newcommand{\milcestacion}[2]{%
\Needspace*{9\baselineskip}%
\stepcounter{milcestacion}%
\pdfbookmark[1]{#2}{milcestacion\themilcestacion}%
\markright{#2}%
\begin{tcolorbox}[enhanced,colback=#1,colframe=#1,arc=2mm,boxrule=0pt,
  left=5mm,right=5mm,top=2.6mm,bottom=2.6mm,before skip=11pt,
  after={\par\nopagebreak\vskip7pt}]
{\sffamily\bfseries\fontsize{15}{18}\selectfont\milctintasobre{#1}#2}
\end{tcolorbox}}

% Tarjeta de paso numerada: sustituye las tablas «Pilar / Aplicacion», que
% eran formato de consulta docente, no de estudiante. El numero va en un
% circulo sobre el borde para que la secuencia se lea de un vistazo.
\newcommand{\milcpaso}[3]{%
\Needspace{4\baselineskip}%
\begin{tcolorbox}[enhanced,colback=white,colframe=milcLinea,arc=1.5mm,boxrule=.9pt,
  left=14mm,right=4mm,top=3mm,bottom=3mm,before skip=4pt,after skip=4pt,
  fontupper=\RaggedRight,
  overlay={\node[anchor=north west,circle,fill=#2,text=white,inner sep=0pt,
    minimum size=7.4mm,font=\sffamily\bfseries\small]
    at ([xshift=3.6mm,yshift=-3.2mm]frame.north west) {#1};}]
#3
\end{tcolorbox}}

% Casilla marcable de tamano util con lapiz (la anterior era \fbox{\phantom{x}},
% de alto variable segun el texto que la seguia).
\newcommand{\milcmarca}{\framebox[3.8mm]{\rule{0pt}{3.4mm}}}

% Fila marcable de lista suelta (checks de la estacion 1).
\newcommand{\milccheckrow}[1]{%
\par\vspace{1.8mm}\noindent
\makebox[6.4mm][l]{\raisebox{-.55mm}{\textcolor{milcNegro}{\milcmarca}}}%
\parbox[t]{\dimexpr\linewidth-6.4mm\relax}{\RaggedRight #1}\par}

% Celda marcable dentro de la rubrica (tabla de autoevaluacion). Misma
% sangria francesa, pero pensada para vivir dentro de una columna X.
\newcommand{\milcopcion}[1]{%
\makebox[5.2mm][l]{\raisebox{-.55mm}{\textcolor{milcNegro}{\milcmarca}}}%
\parbox[t]{\dimexpr\linewidth-5.2mm\relax}{\RaggedRight #1}}

% ── Recursos visuales de la guía ──────────────────────────────────────
% Declarados en el bloque `recursos:` del YAML y emitidos por el builder.
% \guiaFigura   → imagen rasterizada o PDF (foto, carta celeste, montaje).
% \guiaDiagrama → fuente TikZ/circuitikz (.tex) incrustada como vector:
%                 es la vía para circuitos y esquemáticos editables.
% [#1] = texto alternativo (alt) · #2 = ruta relativa al .tex · #3 = pie
% El alt llega al PDF etiquetado (/Alt) y lo lee el lector de pantalla.
\newcommand{\guiaFigura}[3][]{%
\begin{center}
\includegraphics[width=0.88\linewidth,height=0.38\textheight,keepaspectratio,alt={#1}]{#2}\\[1.5mm]
{\footnotesize\itshape\textcolor{milcNegro!75}{#3}}
\end{center}
}

\newcommand{\guiaDiagrama}[2]{%
\begin{center}
\input{#1}\\[1.5mm]
{\footnotesize\itshape\textcolor{milcNegro!75}{#2}}
\end{center}
}

\begin{document}
\thispagestyle{empty}

% ── Portada ───────────────────────────────────────────────────────────
% Antes: un rectangulo de color a pagina completa. Se veia bien en pantalla,
% pero en la fotocopiadora del colegio se comia el toner y salia gris sucio.
% Ahora la banda ocupa el tercio superior y el resto va en blanco.
\begin{tikzpicture}[remember picture,overlay]
  \path[fill=milcGradePrimary]
    (current page.north west) rectangle ([yshift=-10.6cm]current page.north east);

  \node[anchor=north west,text=milcGradeText] at ([xshift=2.0cm,yshift=-1.55cm]current page.north west)
    {\sffamily\bfseries\fontsize{12}{14}\selectfont MOMENTO};
  \node[anchor=north west,text=milcGradeText,opacity=.85] at ([xshift=2.0cm,yshift=-2.35cm]current page.north west)
    {\sffamily\bfseries\fontsize{8.5}{10}\selectfont TERRITORIO INTERIOR · SOCIOEMOCIONAL};
  \node[anchor=north east,text=milcGradeText,opacity=.85,align=right] at ([xshift=-2.0cm,yshift=-1.55cm]current page.north east)
    {\sffamily\bfseries\fontsize{9}{12}\selectfont I.E. Sor María Juliana\\Cartago, Valle del Cauca};

  \node[anchor=west,text=milcGradeText] (gradonum) at ([xshift=1.85cm,yshift=-7.1cm]current page.north west)
    {\sffamily\bfseries\fontsize{112}{112}\selectfont 7};
  % El circulito de grado (°) solo aplica a las guias de grado («11°»); en
  % Bebras y Semillero el numero grande es un momento/modulo, no un grado.
  % Se posiciona relativo al nodo `gradonum`, no en coordenadas absolutas.
  

  \node[anchor=west,text=milcGradeText,text width=11.4cm,align=left] at ([xshift=1.5cm,yshift=0.5cm]gradonum.east)
    {
     {\sffamily\bfseries\fontsize{9}{11}\selectfont\MakeUppercase{Territorio interior · Socioemocional}}\\[3mm]
     {\sffamily\bfseries\fontsize{21}{25}\selectfont\setlength{\baselineskip}{27pt}El que llegó\\[4pt]de otra parte\\[4pt]hasta la ciudad se mudó}\\[3.5mm]
     {\sffamily\bfseries\fontsize{15}{18}\selectfont Grado 9° · Momento 7}
    };
\end{tikzpicture}

\vspace*{9.4cm}
\vfill

{\sffamily\bfseries\fontsize{8}{10}\selectfont\textcolor{milcGradeDark}{LO QUE TE LLEVAS HOY}}

\begin{tcolorbox}[enhanced,colback=white,colframe=milcNegro,arc=2mm,boxrule=1.6pt,
  left=5mm,right=5mm,top=4mm,bottom=4mm,before skip=3pt,after skip=12pt,
  fontupper=\RaggedRight\fontsize{11}{15.5}\selectfont]
Tu mapa de llegada: de dónde vienen las familias de tu curso —contado en generaciones—, lo que necesita de verdad alguien que llega a un colegio nuevo según quienes ya pasaron por eso, y un protocolo de primera semana escrito y propuesto al curso.
\end{tcolorbox}

\begin{tcolorbox}[enhanced,colback=milcGris,colframe=milcGris,arc=2mm,boxrule=0pt,
  left=5mm,right=5mm,top=3.5mm,bottom=3.5mm,after skip=12pt]
{\sffamily\bfseries\fontsize{8}{10}\selectfont\textcolor{milcNegro!55}{ESTA GUÍA ES DE}}\\[3mm]
\begin{tabularx}{\linewidth}{X p{3.2cm} p{3.2cm}}
\linead{\linewidth} & \linead{3.0cm} & \linead{3.0cm}\\[-1mm]
{\fontsize{8.5}{10}\selectfont\textcolor{milcNegro!55}{Nombre completo}} &
{\fontsize{8.5}{10}\selectfont\textcolor{milcNegro!55}{Grupo}} &
{\fontsize{8.5}{10}\selectfont\textcolor{milcNegro!55}{Fecha}}\\
\end{tabularx}
\end{tcolorbox}

\vfill

\noindent\textcolor{milcLinea}{\rule{\linewidth}{1pt}}\\[2mm]
{\sffamily\bfseries\fontsize{9.5}{12}\selectfont Autor: Dr. Álvaro Cárdenas Orozco}\\[1mm]
{\sffamily\fontsize{8.5}{11}\selectfont\textcolor{milcNegro!55}{Metodología MILC · Apertura ancestral · Migración y desarraigo · Triángulo Dussel-Estoicismo-Floridi}}

\newpage

\section*{Apertura: El que llegó de otra parte: hasta la ciudad tuvo que mudarse}

\begin{softbox}{milcGradeDark}{milcGradeSoft}{Saber ancestral}
Antes de hablar de nadie que haya llegado a Cartago, hay que decir algo sobre Cartago: \textbf{esta ciudad llegó de otra parte}. Fue fundada el \textbf{9 de agosto de 1540} por el mariscal \textbf{Jorge Robledo} \emph{a orillas del río Otún} ---es decir, \textbf{en el sitio que hoy ocupa Pereira}---. Y el \textbf{21 de abril de 1691} \textbf{la ciudad entera se trasladó} a donde está ahora, entre el Cauca y La Vieja: unos dicen que por el asedio de los pijaos, otros que por conveniencia agrícola y económica. \emph{Piensa en lo que eso significa, porque es más raro de lo que suena.} \textbf{Un pueblo completo empacó y volvió a empezar en otro sitio} ---y se llevó lo único que no se podía cargar en mula: \textbf{el nombre}---. \emph{Y pasó algo todavía más extraño:} \textbf{siglos después, donde Cartago estuvo, creció otra ciudad}. \emph{Hoy hay dos lugares y una sola historia.} \textbf{Guarda esto, porque desarma de entrada el argumento con que se recibe mal a la gente:} \emph{aquí nadie es de aquí desde siempre. Ni siquiera la ciudad.}
\end{softbox}

\begin{softbox}{milcTurquesa}{milcGris}{Saber contemporáneo}
¿Y cómo es llegar hoy, en concreto? \textbf{Para una parte de quienes llegan a Cartago, llegar es \emph{un trámite}.} Colombia tiene un \textbf{Registro Único de Víctimas} y una \textbf{Unidad para las Víctimas} con oficina en Cartago ---en la \textbf{Casa de Justicia}, barrio San Carlos---, porque este es un municipio \textbf{receptor}: aquí llega gente que tuvo que irse de otro lado. \emph{Y la norma es precisa sobre lo que eso implica.} \textbf{La entidad territorial que recibe a población desplazada debe garantizar} ---mientras se surte la inscripción en el registro--- \textbf{alimentación, artículos de aseo, manejo de abastecimientos, utensilios de cocina y alojamiento transitorio}. \textbf{Y hay un plazo:} quien fue víctima de desplazamiento forzado tiene \textbf{dos años} desde el hecho para pedir su inclusión en el registro. \emph{Detente en la lista de lo que hay que garantizar, porque es un retrato exacto de lo que significa llegar sin nada:} \textbf{no es una metáfora ---son ollas, jabón, comida y un techo por unos días---.} \emph{Y detente en el plazo: llegar, además, tiene fecha de vencimiento para pedir ayuda.}
\end{softbox}

\begin{softbox}{milcMostaza}{milcGris}{Pregunta-puente}
\textit{La ciudad se mudó entera en 1691 \textbf{y se llevó el nombre}. \textbf{¿Cuántas generaciones hace que tu familia está aquí} \ldots \textbf{y de dónde venía antes}?}
\end{softbox}

\begin{softbox}{milcVerde}{milcGris}{Saber y saber hacer}
Al terminar podrás: (1) \textbf{reconocer} que la propia Cartago llegó de otra parte, y usar eso contra la idea de que unos son «de aquí» y otros no; (2) \textbf{describir} lo que la ley exige garantizarle a quien llega desplazado, y qué dice esa lista sobre lo que significa llegar; (3) \textbf{distinguir} lo que se pierde al migrar ---que no es el lugar, es la red---; (4) \textbf{construir} con quienes ya pasaron por eso un protocolo de primera semana para el que llega, y proponerlo.
\end{softbox}



\small
\begin{tabularx}{\textwidth}{>{\bfseries}p{3.0cm}Y}
\rowcolor{milcGradePrimary!12}Autor & Dr. Álvaro Cárdenas Orozco\\
DBA & Comprende los fenómenos que afectan a su territorio, regula sus emociones ante ellos y acompaña a otros con empatía (transversal 6°--11°, articulado con la política «Escuela, territorio de vida» del MEN, Resolución 006519 de 2025, y la Ley 1620 de 2013)\\
Referentes & Primera ayuda psicológica (OMS, 2011/2012) · Directrices IASC (2007) · USGS y Servicio Geológico Colombiano · Pueblo nasa y la minga (CRIC) · Dussel · Estoicismo · Floridi\\
Producto final & Tu mapa de llegada: de dónde vienen las familias de tu curso —contado en generaciones—, lo que necesita de verdad alguien que llega a un colegio nuevo según quienes ya pasaron por eso, y un protocolo de primera semana escrito y propuesto al curso.\\
\end{tabularx}
\normalsize

\puente{En el momento 6 buscaste barreras en el edificio. \textbf{Hoy buscas otra clase de barrera:} \emph{la que encuentra alguien que llega a un lugar donde todo ya está repartido} ---los puestos, los grupos, los apodos, las historias viejas---.}

\milcestacion{milcGradeDark}{Ruta MILC: así vamos a aprender}
\begin{tabularx}{\textwidth}{>{\centering\arraybackslash}X>{\centering\arraybackslash}X>{\centering\arraybackslash}X>{\centering\arraybackslash}X}
\phasecard{1}{milcVerde}{milcGris}{Escucha\\[-1mm]\small Observo, escucho y pregunto.} &
\phasecard{2}{milcTurquesa}{milcGris}{Sistematización\\[-1mm]\small Reconozco lo que cuesta llegar.} &
\phasecard{3}{milcMagenta}{milcGris}{Praxis\\[-1mm]\small Construyo y aplico.} &
\phasecard{4}{milcVino}{milcGris}{Evaluación\\[-1mm]\small Reflexión liberadora.}
\end{tabularx}


\puente{Primero, un mapa que casi nunca se hace: de dónde viene la gente de tu curso, contando generaciones hacia atrás.}

\milcestacion{milcVerde}{Estación 1 de 3 · Escucha}

\begin{softbox}{milcVerde}{milcGris}{Escena para escuchar}
\textbf{Actividad 1 · IDENTIFICA --- ¿Desde cuándo aquí?} \textbf{Primera parte.} Responde por tu propia familia: \emph{¿dónde nacieron tus padres? ¿Y tus abuelos? ¿Y sus padres, si lo sabes?} \textbf{Cuenta cuántas generaciones lleva tu familia en este municipio.} \emph{Si no sabes, escríbelo así ---«no sé»--- y averígualo esta semana: también es un dato.} \textbf{Segunda parte.} En tu curso, \textbf{¿quién llegó de otro lugar?} ---de otro barrio, otro municipio, otro departamento, otro país---. \emph{Sin nombres: cuenta cuántos.} \textbf{Tercera parte.} Si tú llegaste en algún momento a un colegio nuevo, escribe \textbf{cómo fueron tus primeros tres días}: dónde te sentaste, con quién comiste, quién te habló primero, cuánto tardaste en tener a alguien con quien hablar. \emph{Si nunca te ha pasado, escribe qué crees que sentiría alguien} ---y márcalo como lo que es: una suposición---. \textbf{Y la última:} ¿te acuerdas del nombre del último estudiante que llegó a tu curso?
\end{softbox}

{\sffamily\bfseries\fontsize{8}{10}\selectfont\textcolor{milcNegro!55}{MARCA CUANDO LO HAYAS HECHO}}
\milccheckrow{Contaste \textbf{cuántas generaciones} lleva tu familia aquí, o escribiste que no lo sabes.}
\milccheckrow{Contaste cuántos en tu curso llegaron de otro lugar.}
\milccheckrow{Escribiste tus primeros tres días ---o marcaste tu texto como suposición---.}

\begin{infoband}{milcVerde}


\textbf{En tu cuaderno (Actividad 1 · IDENTIFICA).} Título: «Actividad 1. ¿Desde cuándo aquí?». \textbf{Formato:} tus generaciones + cuántos llegaron de otra parte + tus primeros tres días + si recuerdas el nombre del último que llegó. \textbf{Extensión:} media página. \textbf{Sabes que terminaste cuando} sabes \textbf{cuántas generaciones} lleva tu familia en este municipio. Esta página es tuya: \textbf{nadie más tiene que leerla si tú no quieres}.
\end{infoband}

\puente{Ya lo tienes. Ahora, qué se pierde de verdad al migrar y qué exige la ley para quien llega desplazado.}

\milcestacion{milcTurquesa}{Fase 2 · Sistematización: entender la llegada}

\begin{softbox}{milcTurquesa}{milcGris}{Cuatro claves sobre llegar}
Cuatro claves para entender qué cuesta llegar, y por qué el que llega casi nunca lo cuenta.
\end{softbox}

\milcpaso{1}{milcTurquesa}{\textbf{Aquí todos llegamos; solo cambia hace cuánto.} Cartago se fundó en \textbf{1540} donde hoy está Pereira y se trasladó en \textbf{1691}. \emph{El Norte del Valle se pobló después con oleadas sucesivas ---gente que bajó, gente que subió, gente que huyó---, y tu propio mapa de generaciones lo va a mostrar.} \textbf{Por eso «es que ese no es de aquí» es, casi siempre, una frase dicha por alguien que tampoco lo era hace tres generaciones.} \emph{Y funciona igual que la cerca del momento 8 del grado 8:} \textbf{traza una raya en un punto arbitrario del tiempo} ---justo después de que llegó el que la traza---.}
\milcpaso{2}{milcTurquesa}{\textbf{Lo que se pierde al migrar no es el lugar: es la red.} \emph{Uno cree que lo duro de irse es extrañar el paisaje, y no es eso ---o no es lo primero---.} \textbf{Lo primero que se pierde es \emph{la gente que sabía quién eras}:} el que te fiaba, la vecina que te cuidaba, el profesor que sabía que eras bueno en algo, los amigos que conocían tus chistes. \emph{Llegar a un lugar nuevo es empezar sin historia ---nadie sabe nada bueno de ti todavía---.} \textbf{Y ahí hay una injusticia concreta que conviene ver:} \emph{el que llega tiene que \textbf{construir toda su reputación desde cero}, mientras los demás ya tienen la suya hecha}. \textbf{Junta eso con el momento 4 del grado 8 y verás por qué un rumor sobre alguien recién llegado pega tan rápido:} \emph{no hay nada previo que lo contradiga.}}
\milcpaso{3}{milcTurquesa}{\textbf{Llegar, para muchos, es además un trámite.} La norma obliga a la entidad que recibe a garantizar \textbf{alimentación, artículos de aseo, manejo de abastecimientos, utensilios de cocina y alojamiento transitorio} mientras se surte la inscripción en el \textbf{Registro Único de Víctimas}; y hay un plazo de \textbf{dos años} para pedir la inclusión. \emph{Lee esa lista otra vez, porque describe una situación muy concreta:} \textbf{alguien que llegó sin ollas.} \textbf{Y piensa en lo que eso implica para un compañero tuyo:} \emph{papeleo, filas, esperar respuestas, cambiar de casa varias veces} ---mientras se supone que además rinde en clase---. \textbf{Cuando alguien recién llegado falta, llega tarde o no trajo los materiales, esa suele ser la explicación,} \emph{y casi nunca la va a contar.}}
\milcpaso{4}{milcTurquesa}{\textbf{La red no se pide: se ofrece.} \emph{Aquí está el punto práctico de toda la guía.} \textbf{Al que llega no se le puede pedir que se integre}, porque para eso tendría que hacer justo lo que es más difícil sin red: acercarse a un grupo que ya está armado, sin saber sus códigos y arriesgándose a un no. \emph{Es exactamente la situación del momento 9 del grado 7 ---estar por fuera--- pero de estreno.} \textbf{Por eso funciona lo contrario:} \emph{que alguien de adentro haga el primer movimiento, con una invitación concreta} ---qué, cuándo, dónde y una tarea---. \textbf{Y por eso este momento pide un \emph{protocolo} y no buenas intenciones:} \emph{lo que depende de que a alguien se le ocurra, no ocurre.}}

\begin{drawbox}{milcTurquesa}{Lo que de verdad necesita alguien en su primera semana}
\begin{pasoslista}
\item \textbf{Que alguien le diga su nombre y le pregunte el suyo} el primer día. \emph{Suena mínimo; es lo que más recuerdan.}
\item \textbf{Con quién sentarse en el descanso.} \emph{El descanso es la hora más larga del que llega.}
\item \textbf{Las reglas no escritas:} qué profesor es estricto, dónde se compra qué, qué no se hace aquí.
\item \textbf{Estar en el grupo de chat} donde se avisan las cosas. \emph{Faltar ahí es quedar por fuera y además quedar mal.}
\item \textbf{Una tarea:} entrar haciendo algo es infinitamente más fácil que entrar existiendo.
\end{pasoslista}
\end{drawbox}

\begin{softbox}{milcMostaza}{milcGris}{Errores comunes y su corrección}
\textbf{«Que se integre si quiere»:} le pide al que menos red tiene el movimiento más difícil. \textbf{Suponer lo que necesita:} pregúntale a quien ya pasó por eso. \textbf{Presentarlo delante de todos y ya:} eso dura treinta segundos. \textbf{Adoptarlo como proyecto:} se nota y humilla. \textbf{Preguntar por qué se vino:} puede ser lo más difícil de su vida; que lo cuente si quiere. \textbf{Olvidar el grupo de chat:} es la exclusión más fácil de evitar. \textbf{Creer que a la semana ya está:} el momento difícil suele ser el mes dos, cuando dejó de ser novedad.
\end{softbox}

\begin{infoband}{milcTurquesa}


\textbf{En tu cuaderno (Actividad 2 · EXPLICA).} Título: «Actividad 2. Lo que se pierde al llegar». \textbf{Formato:} explica con tus palabras por qué \textbf{lo primero que se pierde es la red y no el lugar}, y por qué eso hace que un rumor sobre alguien recién llegado sea más difícil de desmentir. \textbf{Extensión:} media página. \textbf{Sabes que terminaste cuando} puedes explicar por qué \textbf{«que se integre si quiere» no funciona}.
\end{infoband}

\puente{Ya sabes qué cuesta. Es hora de preguntarle a quien llegó ---que es la única fuente que sirve--- y de escribir un protocolo.}

\milcestacion{milcMagenta}{Fase 3 · Praxis: el protocolo de llegada}

\begin{softbox}{milcMagenta}{milcGris}{Cómo se recibe a alguien que llega}
La única fuente válida aquí son los que ya pasaron por esto. Lo demás es suponer, y suponer es exactamente lo que este grado está aprendiendo a no hacer.
\end{softbox}

\milcpaso{1}{milcMagenta}{\textbf{El mapa de generaciones (10 min, todo el curso).} Junten los datos ---\textbf{sin nombres}--- y hagan el mapa del curso: \emph{cuántas familias llevan aquí más de tres generaciones, cuántas dos, cuántas llegaron con esta generación, de qué lugares}. \textbf{Y averigüen un dato que no está en esta guía:} \emph{cuántas víctimas de desplazamiento hay registradas en Cartago}. \textbf{Búsquenlo en la Unidad para las Víctimas} ---tiene oficina aquí, en la Casa de Justicia--- o en su página. \emph{Este dato no se los damos a propósito: buscarlo es parte de la actividad.}}
\milcpaso{2}{milcMagenta}{\textbf{Preguntarle al que llegó (12 min).} \textbf{Los que llegaron de otro lugar cuentan}, si quieren, \textbf{qué fue lo más difícil de la primera semana} y \textbf{qué habría ayudado}. \emph{Los demás solo escuchan y anotan ---sin arreglar, sin dar consejos: el momento 1 del grado 7---.} \textbf{Y si nadie del curso llegó de otra parte, vayan a otro curso a preguntar.} \emph{Lo que no se puede es inventarlo.}}
\milcpaso{3}{milcMagenta}{\textbf{El protocolo de la primera semana (12 min).} Con lo que dijeron ellos ---\textbf{no con lo que ustedes creen}---, escriban un protocolo: \emph{qué se hace el día uno, qué el día dos, qué la primera semana, qué al mes}. \textbf{Con dos condiciones:} que cada punto tenga \textbf{un responsable con nombre} ---no «el curso»--- y que incluya \textbf{el mes dos}, cuando la novedad ya pasó y es cuando de verdad se decide si la persona se quedó por fuera.}
\milcpaso{4}{milcMagenta}{\textbf{Proponerlo al curso (fuera de clase).} \textbf{Preséntenlo al director de grupo} y pidan que se aplique con el próximo que llegue. \emph{Y háganlo concreto:} \textbf{que quede escrito quién recibe al siguiente}, con nombre, desde ahora. \emph{Un protocolo sin nombre asignado es una intención con formato de documento.}}

\begin{drawbox}{milcMagenta}{Antes de cerrar tu mapa de llegada}
\begin{checklista}
\item El mapa de generaciones está hecho con datos de todo el curso, sin nombres.
\item \textbf{Buscaron ustedes} la cifra de víctimas registradas en Cartago y la anotaron con su fuente.
\item El protocolo salió de lo que dijeron \textbf{quienes llegaron}, no de lo que ustedes suponen.
\item Cada punto tiene un \textbf{responsable con nombre}.
\item El protocolo incluye \textbf{el mes dos}, no solo la primera semana.
\item Ya está escrito \textbf{quién recibe al próximo que llegue}.
\end{checklista}
\end{drawbox}

\begin{softbox}{milcMagenta}{milcGris}{Plantilla de trabajo}
\textbf{Mi mapa.} \emph{«Mi familia lleva \_\_\_\_ generaciones aquí. Antes venía de \_\_\_\_.»} \\[1mm]
\textbf{Para preguntar a quien llegó.} \emph{«¿Qué fue lo más difícil de tu primera semana? ¿Qué te habría ayudado?»} ---escuchar sin arreglar---. \\[1mm]
\textbf{El protocolo.} \emph{Día 1: \_\_\_\_ (responsable: \_\_\_\_) · Semana 1: \_\_\_\_ · Mes 2: \_\_\_\_}
\end{softbox}

\begin{infoband}{milcMagenta}


\textbf{En tu cuaderno (Actividad 3 · CREA).} Título: «Actividad 3. Mi mapa de llegada». \textbf{Formato:} el mapa de generaciones del curso + la cifra que buscaron con su fuente + lo que dijeron quienes llegaron + el protocolo con responsables + a quién se propuso. \textbf{Extensión:} una página. \textbf{Sabes que terminaste cuando} el protocolo tiene \textbf{nombres propios} al frente de cada punto.
\end{infoband}

\puente{El producto es un protocolo con nombres asignados. \emph{«Ser amables con los nuevos» no es un protocolo: no lo hace nadie.}}

\milcestacion{milcMostaza}{Producto de la sesión}
\textbf{Mapa de llegada: generaciones del curso, lo que dijeron quienes llegaron y un protocolo de primera semana con responsables}.
\begin{softbox}{milcMostaza}{milcGris}{Criterios de calidad}
(1) El mapa de generaciones se construye con datos reales del curso, sin identificar familias. (2) La cifra de víctimas registradas en el municipio fue buscada por el grupo y citada con su fuente. (3) El protocolo se basa en lo que relataron quienes llegaron, no en suposiciones del resto. (4) Cada punto del protocolo tiene un responsable con nombre propio. (5) El protocolo contempla el mes dos, cuando la novedad ya pasó.
\end{softbox}

\puente{Antes de cerrar, revisa lo que hace inútil este ejercicio: si tu protocolo salió de lo que tú crees, y no de lo que dijeron los que llegaron, empieza otra vez.}

\milcestacion{milcVino}{Evaluación liberadora · cinco dimensiones}

\begin{softbox}{milcVino}{milcGris}{Sentido de la evaluación}
Evaluar no es solo revisar una nota. Es reconocer qué aprendí, qué cambió en mí, qué emociones manejé, cómo me relaciono con la ciudad y la comunidad, y qué le dejaré a quien viene detrás.
\end{softbox}

\textbf{Mi reflexión liberadora en cinco dimensiones.} Completa cada frase iniciada:

\small
\begin{tabularx}{\textwidth}{>{\bfseries\RaggedRight\arraybackslash}p{3.4cm}Y>{\RaggedRight\arraybackslash}p{4.6cm}}
\rowcolor{milcVino}\color{white}Dimensión & \color{white}\bfseries Mi respuesta (frame starter) & \color{white}\bfseries Algo que descubrí\\
1. Personal & Lo que más cambió en mí fue \linead{6.5cm} & \linead{4.2cm}\\
2. Emocional & La emoción más fuerte fue \linead{2.5cm} y la regulé \linead{3cm} & \linead{4.2cm}\\
3. Ciudadana & En democracia esto importa porque \linead{6cm} & \linead{4.2cm}\\
4. Local (Cartago) & En mi barrio sería útil para \linead{6.5cm} & \linead{4.2cm}\\
5. Intergeneracional & A mi abuela/abuelo le diría \linead{6.5cm} & \linead{4.2cm}\\
\end{tabularx}
\normalsize

\begin{infoband}{milcVino}
\textbf{Para la próxima sustentación.} Vuelve a esta tabla en seis meses. Verás que las respuestas cambian — no porque lo aprendido cambie, sino porque tú cambias. La evaluación liberadora es una conversación contigo a través del tiempo.
\end{infoband}


\puente{Cerramos con tres voces: el que está a la vera del camino reclamando un derecho, el cambio permanente de las cosas y el entorno compartido.}

\milcestacion{milcGradeDark}{Triángulo de pensamiento · cierre filosófico}

\begin{softbox}{milcVino}{milcGris}{Dussel · lente del nosotros}
\textit{«El otro se revela realmente como otro\ldots como el pobre, el oprimido; el que, a la vera del camino, fuera del sistema, muestra su rostro sufriente y sin embargo desafiante: «¡Tengo hambre!, ¡tengo derecho a comer!»»}

{\footnotesize\itshape\color{milcNegro!70}--- Enrique Dussel · Filosofía de la liberación (1977)}

\emph{«A la vera del camino»} es, literalmente, la posición de quien va llegando: \textbf{al lado de un camino que otros ya recorrieron}. \emph{Y la última parte de la cita es exactamente lo que la ley colombiana reconoció:} \textbf{quien llega desplazado no pide caridad ---le corresponden alimentación, aseo, utensilios y alojamiento mientras se resuelve su registro---.} \textbf{Eso convierte la ayuda en otra cosa:} \emph{no es generosidad de quien recibe, es una obligación de quien recibe}. \textbf{Y lo mismo vale, en pequeño, en un curso:} \emph{el puesto del que llega no es un favor que se le hace ---el curso también es suyo desde el primer día---.}

\textbf{Al último que llegó a mi curso, ¿lo tratamos como invitado o como parte?}
\end{softbox}
\linead{\linewidth}\\[1mm]
\linead{\linewidth}

\begin{softbox}{milcMostaza}{milcGris}{Estoicismo · Marco Aurelio · Meditaciones VII, 47 (c. 175 d.C.) · cuidado interior}
\textit{«Contempla las evoluciones de los astros y piensa al mismo tiempo que evolucionas con ellos; y reflexiona continuamente cómo los elementos cambian unos en otros.»}

{\footnotesize\itshape\color{milcNegro!70}--- Marco Aurelio · Meditaciones VII, 47 (c. 175 d.C.)}

\textbf{«Cómo los elementos cambian unos en otros»} describe bastante bien la historia de este territorio: \emph{una ciudad que se llamaba Cartago hoy se llama Pereira, y Cartago está en otro sitio con el mismo nombre}. \textbf{Nada de esto estuvo quieto nunca.} \emph{Y eso sirve para dos cosas.} \textbf{Para el que ya estaba:} \emph{tu familia también llegó ---solo que hace más tiempo---.} \textbf{Y para el que llegó, que es lo que más importa hoy:} \emph{esto también va a cambiar}. \textbf{La primera semana en un colegio nuevo se siente permanente y dura una semana.} \emph{Decirlo no consuela a nadie por sí solo; por eso esta guía termina en un protocolo y no en una frase.}

\textbf{Si mi familia se tuviera que mudar mañana, ¿qué me gustaría que hiciera el curso que me reciba?}
\end{softbox}
\linead{\linewidth}\\[1mm]
\linead{\linewidth}

\begin{softbox}{milcTurquesa}{milcGris}{Floridi · lente de la infoesfera}
\textit{«Somos organismos informacionales ---inforgs--- que compartimos un entorno global hecho, en última instancia, de información: la infoesfera.»}

{\footnotesize\itshape\color{milcNegro!70}--- Luciano Floridi · La cuarta revolución (2014)}

\textbf{Hoy el que llega arrastra algo que antes no existía: su vida anterior sigue disponible.} \emph{Eso tiene una cara buena ---no pierde del todo la red que dejó, puede seguir hablando con sus amigos--- y una difícil:} \textbf{puede quedarse viviendo allá mientras su cuerpo está aquí}, y entonces el mes dos es peor que el uno. \textbf{Y hay una cosa muy concreta que ustedes controlan:} \emph{el grupo de chat}. \textbf{Meter al que llegó el primer día es el gesto de acogida más barato que existe} ---no cuesta nada, no da pena, no hay que hablarle---, \emph{y no hacerlo lo deja por fuera de todo lo que se avisa}. \textbf{Es, de toda esta guía, lo que se puede hacer hoy mismo.}

\textbf{¿El último que llegó a mi curso está en los grupos donde se avisan las cosas?}
\end{softbox}
\linead{\linewidth}\\[1mm]
\linead{\linewidth}

\begin{infoband}{milcNegro}
\textbf{Nota para el docente.} Las tres son citas textuales y llevan su referencia debajo. Puedes remitir a la obra en clase.
\end{infoband}

\milcestacion{milcVino}{Mi autoevaluación · marca una casilla por fila}

{\small
\setlength{\extrarowheight}{2.2pt}
\begin{tabularx}{\textwidth}{>{\RaggedRight\arraybackslash\bfseries}p{3.0cm}YYY}
\rowcolor{milcVino}\color{white}\bfseries Criterio & \color{white}\bfseries Lo logré & \color{white}\bfseries En proceso & \color{white}\bfseries Necesito apoyo\\[.6mm]
Comprensión del tema & \milcopcion{Explico las ideas con mis palabras.} & \milcopcion{Reconozco las ideas, falta precisión.} & \milcopcion{Necesito repasar lo básico.}\\[.8mm]
\rowcolor{milcGris}Calidad del protocolo & \milcopcion{Mi protocolo sale de lo que dijeron quienes llegaron, y tiene tareas asignadas con nombre.} & \milcopcion{Escribí lo que yo creo que necesita alguien nuevo, sin preguntarle.} & \milcopcion{Sigo pensando que el que llega se acomoda solo si quiere.}\\[.8mm]
Aplicación y producto & \milcopcion{Mi producto cumple los criterios.} & \milcopcion{Está armado, con ajustes pendientes.} & \milcopcion{Aún está en construcción.}\\[.8mm]
\rowcolor{milcGris}Reflexión 5 dimensiones & \milcopcion{Respondí cada una con honestidad.} & \milcopcion{Respondí algunas con detalle.} & \milcopcion{Mis respuestas son muy generales.}\\[.8mm]
Triángulo de pensamiento & \milcopcion{Conecté las 3 voces con mi trabajo.} & \milcopcion{Conecté al menos una.} & \milcopcion{No las relaciono con mi proceso.}\\[.8mm]
\end{tabularx}}

\vspace{3mm}
\textbf{Hoy entendí que:}\\[1mm]
\linead{\linewidth}\\[1mm]
\linead{\linewidth}

\textbf{Mi compromiso para la próxima sesión es:}\\[1mm]
\linead{\linewidth}\\[1mm]
\linead{\linewidth}

\begin{softbox}{milcMostaza}{milcGris}{Cita de despedida · Marco Aurelio}
\textit{``El obstáculo en el camino se vuelve el camino. Lo que estorba al camino se vuelve camino.''}\\[1mm]
Cada error de hoy, bien observado, se convierte en pieza del camino que pisarás mañana.
\end{softbox}

\small
\begin{tabularx}{\textwidth}{>{\bfseries}p{3.6cm}Y}
\rowcolor{milcGradeDark!12}Nombre completo: & \linead{10cm}\\
Fecha: & \linead{4cm} \quad Curso/grupo: \linead{4cm}\\
Mi firma: & \linead{9cm}\\
\end{tabularx}
\normalsize

\begin{infoband}{milcGradeDark}
\textbf{Para la próxima sesión.} Dos cosas. \textbf{Primero, traigan el protocolo con los nombres asignados} y \textbf{la cifra que buscaron}, con la fuente de donde la sacaron. \emph{Y cada uno haga hoy mismo lo más barato de la lista: revisar si el último que llegó está en los grupos de chat del curso ---y si no, meterlo---.} \textbf{Segundo, pregunta en tu casa por la llegada:} averigua con alguien mayor \textbf{cómo llegó tu familia a este municipio} ---quién llegó primero, de dónde, por qué se vino, con qué llegó, quién los recibió y cuánto tardaron en sentirse de aquí---. \emph{Si alguien de tu familia tuvo que irse de un lugar por la fuerza, pregúntalo solo si sabes que se puede hablar de eso en tu casa.} \textbf{Trae:} el protocolo y la historia de llegada. \emph{Vas a encontrar casi siempre lo mismo, y por eso vale la pena hacerlo: que alguien los recibió. Un pariente, un paisano, alguien del mismo pueblo que ya estaba aquí. Nadie llega solo del todo, y eso que hizo esa persona es justo lo que ustedes acaban de escribir en un protocolo.}
\end{infoband}

\end{document}
